\documentclass{article}
\usepackage{amsmath}
\usepackage{geometry}
\geometry{a4paper, margin=1in}

\begin{document}

\section*{In-Place FFT Algorithm (Iterative Implementation)}

Given a polynomial $A(x) = a_0 + a_1x + a_2x^2 + \cdots + a_{n-1}x^{n-1}$ with $n = 2^k$ coefficients.

\subsection*{Step 1: Bit-Reversal Permutation}

Rearrange the coefficient array in bit-reversed order. For each index $i$, compute its $k$-bit reversal $\text{rev}(i)$ and swap $a_i$ with $a_{\text{rev}(i)}$ (avoiding duplicate swaps when $i < \text{rev}(i)$).

For example, with $n=8$ (3 bits):
\begin{center}
\begin{tabular}{c|c|c}
Original Index \& Binary \& Bit-Reversed Index \\
\hline
0 \& 000 \& 0 \\
1 \& 001 \& 4 \\
2 \& 010 \& 2 \\
3 \& 011 \& 6 \\
4 \& 100 \& 1 \\
5 \& 101 \& 5 \\
6 \& 110 \& 3 \\
7 \& 111 \& 7 \\
\end{tabular}
\end{center}

The permuted array is:
$$[a_0, a_4, a_2, a_6, a_1, a_5, a_3, a_7]$$

\subsection*{Step 2: Iterative Butterfly Operations}

Starting with $m=2$, double $m$ each iteration until $m=n$. For each stage, we use twiddle factors:
$$\omega_m^j = e^{-2\pi i j / m} = \cos\left(\frac{2\pi j}{m}\right) - i\sin\left(\frac{2\pi j}{m}\right)$$

\subsubsection*{Stage 1: $m=2$ (groups of size 2)}

Process adjacent pairs with butterfly operations:
\begin{align*}
\ &[a_0 + a_4, \; a_0 - a_4, \; a_2 + a_6, \; a_2 - a_6, \; a_1 + a_5, \; a_1 - a_5, \; a_3 + a_7, \; a_3 - a_7] \\
\ &= [A_0^{(1)}, A_4^{(1)}, A_2^{(1)}, A_6^{(1)}, A_1^{(1)}, A_5^{(1)}, A_3^{(1)}, A_7^{(1)}]
\end{align*}

Note: $\omega_2^0 = 1$, so this stage uses only additions and subtractions.

\subsubsection*{Stage 2: $m=4$ (groups of size 4)}

Twiddle factors: $\omega_4^0 = 1$, $\omega_4^1 = -i$

For the first half $[A_0^{(1)}, A_4^{(1)}, A_2^{(1)}, A_6^{(1)}]$:
\begin{align*}
A_0^{(2)} &= A_0^{(1)} + \omega_4^0 \cdot A_2^{(1)} \\
A_2^{(2)} &= A_0^{(1)} - \omega_4^0 \cdot A_2^{(1)} \\
A_4^{(2)} &= A_4^{(1)} + \omega_4^1 \cdot A_6^{(1)} \\
A_6^{(2)} &= A_4^{(1)} - \omega_4^1 \cdot A_6^{(1)}
\end{align*}

For the second half $[A_1^{(1)}, A_5^{(1)}, A_3^{(1)}, A_7^{(1)}]$:
\begin{align*}
A_1^{(2)} &= A_1^{(1)} + \omega_4^0 \cdot A_3^{(1)} \\
A_3^{(2)} &= A_1^{(1)} - \omega_4^0 \cdot A_3^{(1)} \\
A_5^{(2)} &= A_5^{(1)} + \omega_4^1 \cdot A_7^{(1)} \\
A_7^{(2)} &= A_5^{(1)} - \omega_4^1 \cdot A_7^{(1)}
\end{align*}

\subsubsection*{Stage 3: $m=8$ (groups of size 8, final stage)}

Twiddle factors: $\omega_8^0 = 1$, $\omega_8^1 = \frac{\sqrt{2}}{2} - i\frac{\sqrt{2}}{2}$, $\omega_8^2 = -i$, $\omega_8^3 = -\frac{\sqrt{2}}{2} - i\frac{\sqrt{2}}{2}$

For all 8 elements:
\begin{align*}
A_0^{(3)} &= A_0^{(2)} + \omega_8^0 \cdot A_1^{(2)} \\
A_1^{(3)} &= A_0^{(2)} - \omega_8^0 \cdot A_1^{(2)} \\
A_2^{(3)} &= A_2^{(2)} + \omega_8^1 \cdot A_3^{(2)} \\
A_3^{(3)} &= A_2^{(2)} - \omega_8^1 \cdot A_3^{(2)} \\
A_4^{(3)} &= A_4^{(2)} + \omega_8^2 \cdot A_5^{(2)} \\
A_5^{(3)} &= A_4^{(2)} - \omega_8^2 \cdot A_5^{(2)} \\
A_6^{(3)} &= A_6^{(2)} + \omega_8^3 \cdot A_7^{(2)} \\
A_7^{(3)} &= A_6^{(2)} - \omega_8^3 \cdot A_7^{(2)}
\end{align*}

The final result $[A_0^{(3)}, A_1^{(3)}, A_2^{(3)}, A_3^{(3)}, A_4^{(3)}, A_5^{(3)}, A_6^{(3)}, A_7^{(3)}]$ is the DFT.

\subsection*{General In-Place FFT Pseudocode}

\begin{verbatim}
FFT(a, n):
    // Step 1: Bit-reversal permutation
    for i = 0 to n-1:
        j = bit_reverse(i, log2(n))
        if i < j:
            swap(a[i], a[j])

    // Step 2: Iterative butterfly operations
    for m = 2; m <= n; m *= 2:        // m = current group size
        omega_m = exp(-2*pi*i / m)    // principal twiddle factor
        for k = 0; k < n; k += m:     // iterate over groups
            omega = 1                  // omega_m^0
            for j = 0; j < m/2; j++:  // butterfly pairs within group
                u = a[k + j]           // upper butterfly branch
                v = omega * a[k + j + m/2]  // lower butterfly branch
                a[k + j] = u + v       // in-place update
                a[k + j + m/2] = u - v
                omega = omega * omega_m    // update to next twiddle factor
\end{verbatim}

\subsection*{Key Points}

\begin{enumerate}
    \item \textbf{Bit-reversal}: Ensures the input is ordered so the algorithm can run in-place.
    \item \textbf{In-place updates}: Butterfly operations use the same positions for input and output, saving space.
    \item \textbf{Twiddle factor reuse}: Generate $\omega_m^j$ through iterative multiplication $\omega \cdot \omega_m$, avoiding repeated trigonometric calculations.
    \item \textbf{Butterfly structure}: Each pair of elements is computed independently, suitable for parallelization.
\end{enumerate}

\end{document}
